\chapter{Эта-фаза и WZW-ориентация обменной Real-пары}

Предыдущая глава отделила положительный модуль вакуумного отклика от
ориентированного суперследа. Теперь проверяется последняя топологическая
возможность выбора материального сектора: сохраняет ли обменная Real-пара
физическую фазу, отличающую класс пятнадцать от нулевого класса.

\section{Три разных фазовых объекта}

Необходимо различать три конструкции:
\begin{enumerate}
 \item winding петли $V_{15}:S^1\to U(105)$;
 \item нечётный характер Черна её Bott-продолжения на $S^3$;
 \item фаза фермионного детерминанта или Pfaffian line.
\end{enumerate}
Они связаны индексными теоремами, но не являются одной и той же формулой.
Теорема Дая--Чжана связывает Toeplitz index, эта-инвариант и функционал
Весса--Зумино; см. \path{arXiv:math/0103230} и
\path{arXiv:1205.0562}. Экспоненцированный эта-инвариант естественно живёт
в determinant line; см. \path{arXiv:hep-th/9405012}.

Для явной пары
\begin{equation}
 V_{15}(z)=\left(
 zq_0+1-q_0,
 z^{-1}\overline q_0+1-\overline q_0
 \right)
\end{equation}
имеем
\begin{equation}
 \operatorname{wind}V_+=15,
 \qquad
 \operatorname{wind}V_-=-15.
 \label{eq:v5-real-pair-windings}
\end{equation}

\section{Степенное ограничение WZW-формы}

На одномерной окружности кубическая форма тождественно равна нулю:
\begin{equation}
 \Tr\bigl((V^{-1}dV)^3\bigr)=0,
 \qquad \dim S^1=1.
\end{equation}
Поэтому петле нельзя напрямую приписать трёхмерный WZW-функционал. На
$S^1$ её топологический инвариант --- это первая нечётная форма
\begin{equation}
 \frac{1}{2\pi i}\int_{S^1}\Tr(V^{-1}dV)=\pm15.
\end{equation}

После уже построенной относительной Bott-суспензии появляется отображение
$W:S^3\to U$, и его трёхмерный нечётный характер Черна равен
\begin{equation}
 \nu_3(W_+)=15,
 \qquad
 \nu_3(W_-)=-15.
 \label{eq:v5-bott-wzw-charges}
\end{equation}
Следовательно, Real-спаривание обнуляет суммарный ориентированный
WZW-заряд:
\begin{equation}
 \nu_3(W_+)+\nu_3(W_-)=0.
\end{equation}

\section{Комплексная determinant-line фаза}

При целочисленном уровне $k$ экспоненцированная топологическая фаза одной
ветви имеет вид
\begin{equation}
 \exp\bigl(2\pi i k\nu_3(W_\pm)\bigr)=1.
 \label{eq:v5-integer-wzw-phase-trivial}
\end{equation}
Поэтому целочисленный WZW-функционал различает классы до экспоненцирования,
но сам по себе не даёт ненулевой относительной $U(1)$-фазы между сектором
пятнадцать и сектором ноль.

Это согласуется с determinant-line геометрией: изменение reduced
эта-инварианта на целое не меняет его стандартную экспоненту. Для
нетривиальной комплексной фазы потребовался бы дробный уровень, задаваемый
аномалией или дополнительным bulk-объектом. Текущий родитель такого уровня
не выводит.

\section{Условный Pfaffian-знак}

Для одной вещественной ориентированной ветви нечётный спектральный поток
мог бы дать Pfaffian-паритет
\begin{equation}
 \sigma_+=(-1)^{15}=-1.
 \label{eq:v5-reduced-pfaffian-parity}
\end{equation}
Такой знак является физическим только после построения реального
кососамосопряжённого семейства и ориентации его Pfaffian line. Общая связь
знака фермионного интеграла с эта- и mod-two индексами обсуждается в
\path{arXiv:1508.04715} и \path{arXiv:2012.03543}.

В текущей конструкции предъявлен Real-унитарный символ, но не семейство
реальных фермионных операторов с физически ориентированной Pfaffian line.
Поэтому формула \eqref{eq:v5-reduced-pfaffian-parity} является
\emph{условным свидетелем нечётности}, а не уже полученной фазой меры.

Более того, сопряжённая ветвь имеет поток $-15$, но по модулю два
\begin{equation}
 (-1)^{-15}=-1.
\end{equation}
Если обе $J$-сопряжённые ветви входят в полный Pfaffian как независимые
Grassmann-блоки, то
\begin{equation}
 \sigma_{\rm full}=(-1)^{15}(-1)^{-15}=+1.
 \label{eq:v5-full-real-pfaffian-phase}
\end{equation}
Эквивалентно, полный блок имеет структуру
$\mathcal A\oplus(-\overline{\mathcal A})$, и его Pfaffian является
модулем в квадрате. Это точно воспроизводит ранее доказанный гейт
Pfaffian--эта-ориентации версии IV.

\begin{theorem}[Сокращение фазы полной Real-пары]
Целочисленная комплексная WZW-фаза класса пятнадцать тривиальна после
экспоненцирования. Условный reduced Pfaffian-знак одной ориентированной
ветви равен $-1$, но полный обменный KO6-пакет содержит две нечётные
ветви и даёт знак $+1$. Поэтому текущая полная фермионная мера не выбирает
сектор пятнадцать относительно сектора ноль.
\end{theorem}

\section{Сверка с физическим half-count}

Гейт меры версии III доказал, что фермионный интеграл считается
Pfaffian-мерой и не удваивает число физических Dirac-частиц. Это не
означает автоматического сохранения знака одной произвольно выбранной
половины. Half-count относится к модулю и кратностям; ориентация Pfaffian
line является отдельным глобальным данным.

Чтобы сохранить $-1$ одной ветви, необходимо доказать, что физическое
пространство интегрирования является одной reality-орбитой, тогда как
$J$-копия обеспечивает структуру, но не входит вторым независимым
Pfaffian-блоком. Проект такого правила не содержит. Старый
determinant-line inflow gate также показал, что anomaly inflow нельзя
объявить без chiral boundary content, bulk Dirac operator и bordism-класса.

\begin{nogocriterion}[Запрет выбора нечётной половины]
Нельзя использовать $(-1)^{15}$ как механизм выбора материи, одновременно
используя полный KO6-блок для положительного модуля и отбрасывая
сопряжённый Pfaffian только в фазе. Это было бы секторно-зависимым правилом
меры.
\end{nogocriterion}

\section{Статус выбора сектора}

Эта/WZW-путь не уничтожает класс пятнадцать и не ослабляет доказанный
дефект. Он показывает другое: фаза полной Real-меры не объясняет, почему
вакуум должен находиться в ненулевом секторе. Поэтому следующий вопрос
переносится с динамического выбора на глобальное условие носителя:
исключает ли сама хопфово-моритовская архитектура тривиальный класс и
принуждает ли она генератор $c_1=1$ до введения частицы.

\begin{protocol}
\begin{enumerate}
 \item winding ветвей: \textbf{$+15/-15$};
 \item кубическая WZW-форма непосредственно на $S^1$: \textbf{ноль};
 \item трёхмерные Bott-заряды: \textbf{$+15/-15$};
 \item суммарный ориентированный WZW-заряд: \textbf{ноль};
 \item целочисленная экспоненцированная фаза: \textbf{единица};
 \item reduced Pfaffian-паритет: \textbf{условно $-1$};
 \item полный Pfaffian-паритет Real-пары: \textbf{$+1$};
 \item ориентация Pfaffian line: \textbf{не выведена};
 \item выбор сектора пятнадцать фазой: \textbf{нет};
 \item неизбежность дефекта внутри сектора: \textbf{сохраняется};
 \item следующий гейт: \textbf{глобальное принуждение нетривиального сектора}.
\end{enumerate}
\end{protocol}

Машинный сертификат сохранён в
\path{s2t_v5_eta_wzw_real_pair_phase_gate_results.json}.